%!TEX program = xelatex
\documentclass[a4paper,14pt]{extarticle}

% === Кодировка и язык ===
\usepackage{fontspec}
\setmainfont{Liberation Serif}
\setsansfont{Liberation Sans}
\setmonofont{Liberation Mono}

\usepackage{polyglossia}
\setdefaultlanguage{russian}
\setotherlanguage{english}

% === Геометрия страницы ===
\usepackage{geometry}
\geometry{
  a4paper,
  left=30mm, right=15mm,
  top=20mm, bottom=20mm
}

% === Межстрочный интервал ===
\usepackage{setspace}
\onehalfspacing

% === Оглавление и заголовки ===
\usepackage{titlesec}
\titleformat{\section}{\normalfont\large\bfseries}{\thesection.}{0.5em}{}
\titleformat{\subsection}{\normalfont\normalsize\bfseries}{\thesubsection.}{0.5em}{}

% === Колонтитулы ===
\usepackage{fancyhdr}
\setlength{\headheight}{17pt}
\pagestyle{fancy}
\fancyhf{}
\fancyhead[R]{\thepage}
\renewcommand{\headrulewidth}{0pt}

% === Работа с рисунками, таблицами, листингами ===
\usepackage{graphicx}
\usepackage{caption}
\usepackage{listings}
\usepackage{xcolor}
\usepackage{tabularx}
\usepackage{booktabs}
\usepackage{array}
\usepackage{longtable}
\usepackage{float}

% === Листинги кода ===
\lstdefinestyle{Cstyle}{
  language=C,
  basicstyle=\small\ttfamily,
  keywordstyle=\bfseries\color{black},
  commentstyle=\itshape\color{gray},
  stringstyle=\color{black},
  numbers=left,
  numberstyle=\tiny,
  stepnumber=1,
  numbersep=6pt,
  breaklines=true,
  frame=single,
  tabsize=2,
  showstringspaces=false,
  extendedchars=true,
  inputencoding=utf8
}

\lstdefinestyle{ASMstyle}{
  basicstyle=\small\ttfamily,
  numbers=none,
  breaklines=true,
  frame=single,
  tabsize=8,
  showstringspaces=false
}

\lstdefinestyle{PLIstyle}{
  basicstyle=\small\ttfamily,
  numbers=none,
  breaklines=true,
  frame=single,
  showstringspaces=false
}

% === Гиперссылки ===
\usepackage[hidelinks]{hyperref}

% === Прочее ===
\usepackage{enumitem}
\usepackage{amsmath}
\usepackage{indentfirst}
\setlength{\parindent}{1.25cm}

% === Подписи к рисункам ===
\captionsetup[figure]{name=Рис.,labelsep=period}
\captionsetup[table]{name=Таблица,labelsep=period}

\begin{document}

% ============================================================
% ТИТУЛЬНЫЙ ЛИСТ
% ============================================================
\thispagestyle{empty}
\begin{center}
  \textbf{Санкт-Петербургский политехнический университет}\\
  \textbf{Институт компьютерных наук и технологий}\\
  \textbf{Высшая школа программной инженерии}\\[3cm]

  {\Large\textbf{К У Р С О В А Я \quad Р А Б О Т А}}\\[1cm]
  {\large Разработка учебной системы программирования}\\[0.5cm]
  {\large Компилятор с ЯВУ}\\[2cm]

  Дисциплина: «Системы программирования»\\[1cm]
  Вариант~№~11\\[2cm]
\end{center}

\begin{flushright}
  \begin{tabular}{ll}
    Исполнители: & Студенты группы~5140904/50102\\
                 & Сырбу~Г.~В.\\
                 & Федерова~С.~К.\\[0.5cm]
    Руководитель: & Расторгуев~В.~Я.\\
  \end{tabular}
\end{flushright}

\vfill
\begin{center}
  Санкт-Петербург\\
  2026
\end{center}

\newpage

% ============================================================
% СОДЕРЖАНИЕ
% ============================================================
\tableofcontents
\newpage

% ============================================================
% РАЗДЕЛ 1. ЗАДАНИЕ НА КУРСОВУЮ РАБОТУ
% ============================================================
\section{Задание на курсовую работу}

Задание на курсовую работу дано в следующей общей формулировке:

\begin{quote}
«Необходимо выполнить доработку элементов макета учебной системы
программирования до уровня, позволяющего обрабатывать «новые» для макета
конструкции языка высокого уровня, примененные в соответствующем варианте».
\end{quote}

Общая формулировка персонализирована в варианте~№~11, приведённом на рис.~\ref{fig:variant}.

\begin{figure}[H]
\centering
\begin{lstlisting}[style=PLIstyle]
EX11: PROC OPTIONS ( MAIN );
DCL A BIN FIXED ( 15 ) INIT ( 11B );
DCL B DEC FIXED ( 3  ) INIT ( 11B );
DCL C BIN FIXED ( 15 );
C = B * 5;
C = A = C;
END EX11;
\end{lstlisting}
\caption{Персональный вариант~№~11}
\label{fig:variant}
\end{figure}

Схема функционирования макета учебной системы программирования, являющегося
согласно заданию на КР одним из элементов исходных данных, приведена на рис.~\ref{fig:arch}.

\begin{figure}[H]
\centering
\begin{tabular}{|c|c|c|c|c|}
  \hline
  \begin{tabular}[t]{c}Исходный\\текст\\на ЯВУ\end{tabular} &
  \(\rightarrow\) &
  \begin{tabular}[t]{c}Кросс-\\компилятор\\с ЯВУ\end{tabular} &
  \(\rightarrow\) &
  \begin{tabular}[t]{c}Эквивалент\\исх.~текста\\на Ассемблере\end{tabular}\\
  \hline
\end{tabular}

\bigskip
\begin{tabular}{|c|c|c|c|c|}
  \hline
  \begin{tabular}[t]{c}Эквивалент\\на Ассемблере\end{tabular} &
  \(\rightarrow\) &
  \begin{tabular}[t]{c}Кросс-\\компилятор\\с Ассемблера\end{tabular} &
  \(\rightarrow\) &
  \begin{tabular}[t]{c}Объектный\\эквивалент\end{tabular}\\
  \hline
\end{tabular}

\bigskip
\begin{tabular}{|c|c|c|c|c|}
  \hline
  \begin{tabular}[t]{c}Объектный\\эквивалент\end{tabular} &
  \(\rightarrow\) &
  \begin{tabular}[t]{c}Загрузчик /\\Эмулятор /\\Отладчик\end{tabular} &
  \(\rightarrow\) &
  \begin{tabular}[t]{c}Загрузочный\\эквивалент\end{tabular}\\
  \hline
\end{tabular}
\caption{Схема функционирования учебной системы программирования}
\label{fig:arch}
\end{figure}

Дополнительные требования к КР предполагают, что данная система программирования
работает на технологической ЭВМ (IBM PC) и является по существу кросс-системой для
объектной ЭВМ (IBM~370). В этой системе:
\begin{itemize}
  \item в качестве языка высокого уровня (ЯВУ) выбран язык, образованный из
        подмножества языковых конструкций PL/1; исходная программа готовится в виде
        текстового файла с расширением \texttt{*.pli};
  \item язык Ассемблера сформирован из языковых конструкций Ассемблера IBM~370;
        ассемблерный эквивалент исходной программы формируется в виде текстового
        файла с расширением \texttt{*.ass};
  \item объектный эквивалент исходной программы готовится в формате объектных файлов
        ОС IBM~370 и хранится в виде двоичного файла с расширением \texttt{*.tex};
  \item загрузочный эквивалент исходной программы представляет собой машинный код
        IBM~370, запоминаемый в области ОЗУ технологической ЭВМ, являющейся зоной
        загрузки для эмулятора объектной ЭВМ.
\end{itemize}

В данном отчёте документируются результаты модификации макетного компилятора с~ЯВУ
под персональный вариант~№~11.

\subsection*{Контрольный (демонстрационный) пример}

Контрольным примером, на который изначально настроен макет компилятора,
является программа \texttt{examppl.pli} (рис.~\ref{fig:demo}).

\begin{figure}[H]
\centering
\begin{lstlisting}[style=PLIstyle]
EXAMP: PROC OPTIONS ( MAIN );
DCL A BIN FIXED ( 31 ) INIT ( 11B  );
DCL B BIN FIXED ( 31 ) INIT ( 100B );
DCL C BIN FIXED ( 31 ) INIT ( 101B );
DCL D BIN FIXED ( 31 );
D = A + B - C;
END EXAMP;
\end{lstlisting}
\caption{Демонстрационный пример (\texttt{examppl.pli})}
\label{fig:demo}
\end{figure}

\subsection*{Отличия персонального варианта от демонстрационного примера}

Персональный вариант~№~11 приведён ниже:

\begin{lstlisting}[style=PLIstyle]
EX11: PROC OPTIONS ( MAIN );
DCL A BIN FIXED ( 15 ) INIT ( 11B );
DCL B DEC FIXED ( 3  ) INIT ( 11B );
DCL C BIN FIXED ( 15 );
C = B * 5;
C = A = C;
END EX11;
\end{lstlisting}

При визуальном сравнении персонального варианта~№~11 и демонстрационного примера
(рис.~\ref{fig:demo}) выявлены следующие отличия:

\begin{enumerate}
  \item В персональном варианте введена переменная \texttt{B} типа
        \texttt{DEC FIXED} (упакованный десятичный формат), тогда как в
        демонстрационном примере все переменные имеют тип \texttt{BIN FIXED}
        (двоичный формат с фиксированной точкой).

  \item В операторе \texttt{C = B * 5;} применяется операция умножения «\texttt{*}»
        упакованного десятичного операнда на целочисленный литерал (десятичную
        константу), отсутствующая в демонстрационном примере.

  \item В операторе \texttt{C = A = C;} применяется оператор сравнения «\texttt{=}»
        в правой части оператора присваивания, формирующий булево (логическое)
        значение — конструкция, отсутствующая в демонстрационном примере.
\end{enumerate}

% ============================================================
% РАЗДЕЛ 2. ПОСТАНОВКА ЗАДАЧИ
% ============================================================
\section{Постановка задачи модификации макета в части компилятора с ЯВУ}

Схема функционирования модифицированного макетного компилятора с~ЯВУ должна
соответствовать рис.~\ref{fig:compil_scheme}.

\begin{figure}[H]
\centering
\begin{tabular}{|c|c|c|}
  \hline
  \begin{tabular}[t]{c}Исходный текст\\на ЯВУ\\\texttt{(ex11.pli)}\end{tabular} &
  \fbox{\begin{tabular}[t]{c}\textbf{Компилятор}\\с ЯВУ\\\texttt{(komppl.exe)}\end{tabular}} &
  \begin{tabular}[t]{c}Эквивалент\\исх. текста на\\Ассемблере\\\texttt{(ex11.ass)}\end{tabular}\\
  \hline
\end{tabular}
\caption{Схема функционирования модифицированного макетного компилятора с~ЯВУ}
\label{fig:compil_scheme}
\end{figure}

Функционал модифицированного макетного компилятора должен обеспечить компиляцию
исходного текста на ЯВУ (рис.~\ref{fig:variant}) в эквивалент исходного текста на
Ассемблере, представленный на рис.~\ref{fig:etalon}.

\begin{figure}[H]
\centering
\begin{lstlisting}[style=ASMstyle]
  Метка      КОП    Операнды         Пояснения
EX11     START  0            Init rel addr=0
         BALR   @RBASE,0     Set @RBASE=cur PC
         USING  *,@RBASE     Use @RBASE for addr
         ZAP    @BUF(8),B(3) Copy pdec B->BUF
         MP     @BUF(8),@FIVE(1) Mul BUF by 5->BUF
         CVB    @R1,@BUF     pdec BUF ->int @R1
         STH    @R1,C        Store int ->C
         LH     @R2,A        Load A ->@R2
         CH     @R2,C        Compare @R2 with C
         BC     8,@TRUE      Branch if equal
         LA     @R3,0        Set @R3=0 (F)
         B      @SAVE        Jump to save
@TRUE    LA     @R3,1        Set @R3=1 (T)
@SAVE    STH    @R3,C        Store bool ->C
         BCR    15,@RVIX     Return to caller
A        DC     H'3'         Store halfword init
B        DC     PL3'3'       Store pdec init
C        DS     H            Reserve halfword
@FIVE    DC     PL1'5'       Store packed const
         DS     0F           Align next item
@BUF     DC     PL8'0'       Clear packed BUF
@RBASE   EQU    5            Set @RBASE=5
@R1      EQU    6            Set @R1=6
@R2      EQU    2            Set @R2=2
@R3      EQU    3            Set @R3=3
@RVIX    EQU    14           Set @RVIX=14
         END                 Program ends
\end{lstlisting}
\caption{Эталон ассемблерного эквивалента персонального варианта~№~11}
\label{fig:etalon}
\end{figure}

\subsection*{Согласованные функциональные ограничения}

По согласованию с преподавателем в модифицированном макетном компиляторе введены
следующие функциональные ограничения:

\begin{enumerate}
  \item Компилятор не различает приоритеты арифметических операций: операция
        умножения «\texttt{*}» в правой части оператора присваивания
        вычисляется строго последовательно слева направо, без учёта её более
        высокого приоритета относительно операций сложения и вычитания.

  \item Литерал-множитель в операции «\texttt{*}» является десятичной целочисленной
        константой в диапазоне~$0$--$9$ (одна цифра), что позволяет ограничиться
        одним байтом в упакованном десятичном представлении.

  \item Оператор сравнения «\texttt{=}» в правой части присваивания допускает
        только сравнение двух переменных типа \texttt{BIN FIXED} и формирует
        булев результат (0~или~1), сохраняемый в переменную типа \texttt{BIN FIXED}.
\end{enumerate}

% ============================================================
% РАЗДЕЛ 3. МОДИФИКАЦИЯ БД КОМПИЛЯТОРА
% ============================================================
\section{Модификация состава и структуры БД макета компилятора с ЯВУ}

\subsection{Модификация БД лексического анализатора}

БД лексического анализатора определяется путём выделения подмножества терминальных
символов-разделителей термов языка. Функция лексического анализа в макете реализована
в подпрограмме \texttt{compress\_ISXTXT()}.

Для персонального варианта~№~11 разделителями термов являются терминальные символы
из множества:
\[
  \{\,\texttt{":"}\,,\, \texttt{"("}\,,\, \texttt{")"}\,,\, \texttt{";"}\,,\,
     \texttt{"+"}\,,\, \texttt{"-"}\,,\, \texttt{"="}\,,\, \texttt{"\_"}\,,\,
     \texttt{"*"}\,\}
\]

Модификация лексических настроек демонстрационного примера состоит в добавлении
символа умножения «\texttt{*}» в перечень разделителей термов. Данный символ в
оригинальном макете не входил в число обрабатываемых разделителей. Поскольку авторы
макета разместили перечень символов-разделителей непосредственно в теле подпрограммы
\texttt{compress\_ISXTXT()}, модификация выполнена в теле этой подпрограммы:

\begin{lstlisting}[style=Cstyle,caption={Фрагмент \texttt{compress\_ISXTXT()} с добавленным символом «\texttt{*}»},label=lst:lex]
void compress_ISXTXT()
{
  /* ... */
  if (
    ( ISXTXT[I1][I2] == '+' ||
      ISXTXT[I1][I2] == '-' ||
      ISXTXT[I1][I2] == '=' ||
      ISXTXT[I1][I2] == '(' ||
      ISXTXT[I1][I2] == ')' ||
      ISXTXT[I1][I2] == '*'   /* <- добавлен разделитель '*' */
    )
    && PREDSYM == ' '
  )
  { I3--; goto L1; }

  if ( ISXTXT[I1][I2] == ' ' &&
     ( PREDSYM == '+' || PREDSYM == '-' ||
       PREDSYM == '=' || PREDSYM == '*'   /* <- добавлен '*' */
     )
  )
  { goto L2; }
  /* ... */
}
\end{lstlisting}

\subsection{Модификация БД синтаксического анализатора}

Высокоуровневое определение модифицированного синтаксиса для персонального
варианта~№~11 приведено ниже; модификации по сравнению с грамматикой
демонстрационного примера выделены курсивом:

\begin{enumerate}
  \item $\langle\text{PRO}\rangle ::= \langle\text{OPR}\rangle\langle\text{TEL}\rangle\langle\text{OEN}\rangle$
  \item $\langle\text{OPR}\rangle ::= \langle\text{IPR}\rangle\texttt{:PROC\_OPTIONS(MAIN);}$
  \item $\langle\text{IPR}\rangle ::= \langle\text{IDE}\rangle$
  \item $\langle\text{IDE}\rangle ::= \langle\text{BUK}\rangle \mid \langle\text{IDE}\rangle\langle\text{BUK}\rangle \mid \langle\text{IDE}\rangle\langle\text{CIF}\rangle$
  \item $\langle\text{BUK}\rangle ::= \texttt{A} \mid \texttt{B} \mid \texttt{C} \mid \texttt{D} \mid \texttt{E} \mid \texttt{M} \mid \texttt{P} \mid \texttt{X}$
  \item $\langle\text{CIF}\rangle ::= \texttt{0} \mid \texttt{1} \mid \texttt{2} \mid \texttt{3} \mid \texttt{4} \mid \texttt{5} \mid \texttt{6} \mid \texttt{7} \mid \texttt{8} \mid \texttt{9}$
  \item $\langle\text{TEL}\rangle ::= \langle\text{ODC}\rangle \mid \langle\text{TEL}\rangle\langle\text{ODC}\rangle \mid \langle\text{TEL}\rangle\langle\text{OPA}\rangle$
  \item $\langle\text{ODC}\rangle ::= \texttt{DCL\_}\langle\text{IPE}\rangle\texttt{\_BIN\_FIXED(}\langle\text{RZR}\rangle\texttt{);}$\\
        $\mid\; \texttt{DCL\_}\langle\text{IPE}\rangle\texttt{\_BIN\_FIXED(}\langle\text{RZR}\rangle\texttt{)INIT(}\langle\text{LIT}\rangle\texttt{);}$\\
        \textit{$\mid\; \texttt{DCL\_}\langle\text{IPE}\rangle\texttt{\_DEC\_FIXED(}\langle\text{RZR}\rangle\texttt{);}$}\\
        \textit{$\mid\; \texttt{DCL\_}\langle\text{IPE}\rangle\texttt{\_DEC\_FIXED(}\langle\text{RZR}\rangle\texttt{)INIT(}\langle\text{LIT}\rangle\texttt{);}$}
  \item $\langle\text{IPE}\rangle ::= \langle\text{IDE}\rangle$
  \item $\langle\text{RZR}\rangle ::= \langle\text{CIF}\rangle \mid \langle\text{RZR}\rangle\langle\text{CIF}\rangle$
  \item $\langle\text{LIT}\rangle ::= \langle\text{MAN}\rangle\texttt{B}$
  \item $\langle\text{MAN}\rangle ::= \texttt{1} \mid \langle\text{MAN}\rangle\texttt{0} \mid \langle\text{MAN}\rangle\texttt{1}$
  \item $\langle\text{OPA}\rangle ::= \langle\text{IPE}\rangle\texttt{=}\langle\text{AVI}\rangle\texttt{;}$
  \item $\langle\text{AVI}\rangle ::= \langle\text{LIT}\rangle \mid \langle\text{IPE}\rangle \mid $\textit{$\langle\text{RZR}\rangle$}$\mid \langle\text{AVI}\rangle\langle\text{ZNK}\rangle\langle\text{LIT}\rangle$\\
        $\mid\; \langle\text{AVI}\rangle\langle\text{ZNK}\rangle\langle\text{IPE}\rangle$
        \textit{$\mid\; \langle\text{AVI}\rangle\langle\text{ZNK}\rangle\langle\text{RZR}\rangle$}
  \item $\langle\text{ZNK}\rangle ::= \texttt{+} \mid \texttt{-}$ \textit{$\mid\; \texttt{*} \mid \texttt{=}$}
  \item $\langle\text{OEN}\rangle ::= \texttt{END\_}\langle\text{IPR}\rangle$
\end{enumerate}

Пояснения к метасимволам:
\begin{itemize}
  \item \texttt{<} и \texttt{>} — ограничители нетерминального символа;
  \item \texttt{::=} — «равно по определению»;
  \item \texttt{|} — альтернативное определение («или»);
  \item \texttt{\_} — терминальный символ «пробел».
\end{itemize}

Перечень нетерминальных символов приведён в таблице~\ref{tab:nonterm}.

\begin{table}[H]
\centering
\caption{Нетерминальные символы грамматики}
\label{tab:nonterm}
\begin{tabularx}{\textwidth}{|l|X|}
  \hline
  \textbf{Нетерминал} & \textbf{Смысл} \\
  \hline
  \texttt{<PRO>} & программа \\
  \texttt{<OPR>} & оператор пролога программы \\
  \texttt{<IPR>} & имя программы \\
  \texttt{<IDE>} & идентификатор \\
  \texttt{<BUK>} & буква \\
  \texttt{<CIF>} & цифра \\
  \texttt{<TEL>} & тело программы \\
  \texttt{<ODC>} & оператор объявления переменных (DCL) \\
  \texttt{<IPE>} & имя переменной \\
  \texttt{<RZR>} & разрядность / целочисленный литерал \\
  \texttt{<LIT>} & двоичный литерал \\
  \texttt{<MAN>} & мантисса литерала \\
  \texttt{<OPA>} & оператор присваивания \\
  \texttt{<AVI>} & арифметическое выражение \\
  \texttt{<ZNK>} & знак операции \\
  \texttt{<OEN>} & оператор эпилога программы (END) \\
  \hline
\end{tabularx}
\end{table}

По сравнению с грамматикой демонстрационного примера модифицированы правила:
\begin{itemize}
  \item \textbf{Правило~8 (\texttt{<ODC>}):} добавлены два новых альтернативных
        правила для объявления переменных типа \texttt{DEC FIXED} — без
        инициализации и с инициализацией;
  \item \textbf{Правило~14 (\texttt{<AVI>}):} добавлена возможность применения
        нетерминала~\texttt{<RZR>} (целочисленный литерал) как самостоятельного
        операнда выражения, а также как правого операнда в паре
        \texttt{<AVI><ZNK><RZR>};
  \item \textbf{Правило~15 (\texttt{<ZNK>}):} добавлены два новых знака операции:
        «\texttt{*}» (умножение) и «\texttt{=}» (сравнение).
\end{itemize}

Состав нетерминальных символов после модификации не изменился.

\subsubsection*{Модификация таблицы синтаксических правил \texttt{SINT[]}}

Таблица \texttt{SINT[]} реализует двунаправленный список с альтернативными
ветвлениями, в котором каждая запись описывает один шаг распознавания (форма
распознавания — от правой части к нетерминалу). В ходе модификации:

\begin{itemize}
  \item размер таблицы увеличен с~201 до~232 записей (\texttt{\#define NSINT~232});
  \item добавлены записи с номерами 201--231, описывающие новые ветки грамматики.
\end{itemize}

Добавленные записи таблицы \texttt{SINT[]} приведены в листинге~\ref{lst:sint_new}.

\begin{figure}[H]
\begin{lstlisting}[style=Cstyle,caption={Новые записи таблицы \texttt{SINT[]} (правила 201--231)},label=lst:sint_new]
 /*                            DEC FIXED path (alt from 187)               */
 {/*.  201     .*/   202 ,    49 , "D  " ,    0 },
 {/*.  202     .*/   203 ,   201 , "E  " ,    0 },
 {/*.  203     .*/   204 ,   202 , "C  " ,    0 },
 {/*.  204     .*/   205 ,   203 , "   " ,    0 },
 {/*.  205     .*/   206 ,   204 , "F  " ,    0 },
 {/*.  206     .*/   207 ,   205 , "I  " ,    0 },
 {/*.  207     .*/   208 ,   206 , "X  " ,    0 },
 {/*.  208     .*/   209 ,   207 , "E  " ,    0 },
 {/*.  209     .*/   210 ,   208 , "D  " ,    0 },
 {/*.  210     .*/   211 ,   209 , "(  " ,    0 },
 {/*.  211     .*/   212 ,   210 , "RZR" ,    0 },
 {/*.  212     .*/   213 ,   211 , ")  " ,    0 },
 {/*.  213     .*/   214 ,   212 , ";  " ,  216 },/* ALT=216: INIT branch  */
 {/*.  214     .*/   215 ,   213 , "ODC" ,    0 },
 {/*.  215     .*/     0 ,   214 , "*  " ,    0 },
 {/*.  216     .*/   217 ,   212 , "I  " ,    0 },/* DEC FIXED with INIT   */
 {/*.  217     .*/   218 ,   216 , "N  " ,    0 },
 {/*.  218     .*/   219 ,   217 , "I  " ,    0 },
 {/*.  219     .*/   220 ,   218 , "T  " ,    0 },
 {/*.  220     .*/   221 ,   219 , "(  " ,    0 },
 {/*.  221     .*/   222 ,   220 , "LIT" ,    0 },
 {/*.  222     .*/   223 ,   221 , ")  " ,    0 },
 {/*.  223     .*/   224 ,   222 , ";  " ,    0 },
 {/*.  224     .*/   225 ,   223 , "ODC" ,    0 },
 {/*.  225     .*/     0 ,   224 , "*  " ,    0 },
 /*                            = as ZNK (comparison operator)               */
 {/*.  226     .*/   227 ,     0 , "=  " ,    0 },
 {/*.  227     .*/   228 ,   226 , "ZNK" ,    0 },
 {/*.  228     .*/     0 ,   227 , "*  " ,    0 },
 /*                            RZR as AVI operand                           */
 {/*.  229     .*/   230 ,   165 , "RZR" ,    0 },
 {/*.  230     .*/   231 ,   229 , "AVI" ,    0 },
 {/*.  231     .*/     0 ,   230 , "*  " ,    0 }
\end{lstlisting}
\end{figure}

Кроме добавления новых записей, в существующих записях изменены поля альтернативы
\texttt{ALT}:
\begin{itemize}
  \item запись~168: поле \texttt{ALT} изменено с~0 на~229\\
        (добавление ветки \texttt{<AVI><ZNK><RZR>});
  \item запись~187: поле \texttt{ALT} изменено с~0 на~201\\
        (добавление ветки \texttt{DEC~FIXED}).
\end{itemize}

\subsubsection*{Модификация таблицы входов \texttt{VXOD[]} и матрицы \texttt{TPR[][]}}

В таблице входов \texttt{VXOD[]} для терминального символа «\texttt{=}» изменено
поле \texttt{VX}: значение изменено с~0 на~226, что обеспечивает возможность
начала разбора нетерминала \texttt{<ZNK>} при встрече символа «\texttt{=}»:

\begin{lstlisting}[style=Cstyle,caption={Изменение в \texttt{VXOD[]} для символа «=»},label=lst:vxod]
  {/*.  50     .*/   "=  " , 226 , 'T' },
\end{lstlisting}

В матрице предшествования \texttt{TPR[][]} для строки символа «\texttt{=}» в
столбце нетерминала \texttt{ZNK} (индекс~15) значение изменено с~0 на~1:

\begin{lstlisting}[style=Cstyle,caption={Изменение в \texttt{TPR[][]} для символа «=» и нетерминала ZNK},label=lst:tpr]
  {/*  =*/ 0,0,0,0,0,0,0,0,0,0,0,0,0,0,0, 1 },
\end{lstlisting}

% ============================================================
% РАЗДЕЛ 4. МОДИФИКАЦИЯ АЛГОРИТМА
% ============================================================
\section{Модификация алгоритма макета компилятора с ЯВУ}

Модификация алгоритма макета компилятора с~ЯВУ произведена в части семантических
подпрограмм. Состав нетерминальных символов после модификации синтаксиса не
изменился, поэтому задача синтеза новых семантических подпрограмм отсутствует.
Анализ изменений в правилах грамматики позволяет определить перечень подпрограмм,
подлежащих модификации:

\begin{itemize}
  \item \textbf{\texttt{ODC1()}} — модифицирована, так как изменена правая часть
        правила для нетерминала \texttt{<ODC>} (добавлен тип \texttt{DEC FIXED});
  \item \textbf{\texttt{AVI2()}} — модифицирована, так как изменены правила для
        нетерминалов \texttt{<AVI>} и \texttt{<ZNK>} (добавлены операции «\texttt{*}» и «\texttt{=}»,
        а также операнд \texttt{<RZR>});
  \item \textbf{\texttt{OPA2()}} — модифицирована в части использования нового
        механизма передачи результата между подпрограммами;
  \item \textbf{\texttt{OEN2()}} — модифицирована для корректной генерации секции
        данных ассемблерного текста с учётом нового типа \texttt{DEC FIXED};
  \item \textbf{\texttt{OPR2()}} — модифицирована в части именования регистров
        и формирования комментариев.
\end{itemize}

Дополнительно введены два вспомогательных механизма:
\begin{itemize}
  \item глобальные переменные \texttt{RESULT\_REG}, \texttt{PENDING\_LABEL},
        \texttt{EXTRA[]}, \texttt{NEXTRA} для передачи состояния между
        семантическими подпрограммами;
  \item подпрограммы \texttt{SETCOMM()} и \texttt{const\_label()} для
        формирования комментариев и меток констант.
\end{itemize}

\subsection{Новые глобальные переменные}

\begin{lstlisting}[style=Cstyle,caption={Новые глобальные переменные},label=lst:globals]
char RESULT_REG[9]  = "@R2";  /* регистр-результат последнего выражения */
char PENDING_LABEL[9] = "";   /* метка, ожидающая навешивания на следующую
                                 генерируемую карточку (используется при
                                 генерации булева сравнения)              */
struct {
    char name[9];   /* символическое имя константы (@FIVE и т.п.) */
    char type;      /* тип ('P' = packed decimal)                  */
    int  size;      /* размер в байтах                             */
    char value[10]; /* строковое значение                          */
} EXTRA[20];
int NEXTRA = 0;     /* количество записей в EXTRA                  */
\end{lstlisting}

\subsection{Вспомогательная подпрограмма \texttt{SETCOMM()}}

Подпрограмма предназначена для записи текстового комментария в поле \texttt{COMM}
генерируемой ассемблерной карточки с учётом возможного переполнения поля операнда:

\begin{lstlisting}[style=Cstyle,caption={Подпрограмма \texttt{SETCOMM()}},label=lst:setcomm]
#define COMM_LIM 25
void SETCOMM ( const char *s )
{
  int i;
  /* Поле COMM может содержать хвост длинного операнда,
     вышедшего за границу фиксированного поля операндов.
     Сохраняем ненулевой префикс, добавляем разделитель. */
  int prefix = 0;
  while ( prefix < 52 && ASS_CARD._BUFCARD.COMM[prefix] != ' ' )
    prefix++;
  int start = prefix;
  if ( prefix > 0 && prefix + 1 < 52 ) start = prefix + 1;
  for ( i = start; i < 52; i++ )
    ASS_CARD._BUFCARD.COMM[i] = ' ';
  if ( !s ) return;
  for ( i = 0; i < COMM_LIM && s[i] != '\0' && start + i < 52; i++ )
    ASS_CARD._BUFCARD.COMM[start + i] = s[i];
}
\end{lstlisting}

\subsection{Вспомогательная подпрограмма \texttt{const\_label()}}

Подпрограмма формирует символическое имя для десятичной константы
(например, «5» → «\texttt{@FIVE}»):

\begin{lstlisting}[style=Cstyle,caption={Подпрограмма \texttt{const\_label()}},label=lst:constlabel]
void const_label ( char *buf, const char *val )
{
  long v = atol(val);
  switch ( v ) {
    case 0: sprintf(buf, "@ZERO");  break;
    case 1: sprintf(buf, "@ONE");   break;
    case 2: sprintf(buf, "@TWO");   break;
    case 3: sprintf(buf, "@THREE"); break;
    case 4: sprintf(buf, "@FOUR");  break;
    case 5: sprintf(buf, "@FIVE");  break;
    case 6: sprintf(buf, "@SIX");   break;
    case 7: sprintf(buf, "@SEVEN"); break;
    case 8: sprintf(buf, "@EIGHT"); break;
    case 9: sprintf(buf, "@NINE");  break;
    default: sprintf(buf, "@C%s", val); break;
  }
}
\end{lstlisting}

\subsection{Модификация подпрограммы \texttt{ODC1()}}

В подпрограмму первого прохода \texttt{ODC1()} добавлена обработка нового типа
\texttt{DEC FIXED}: при наличии слов \texttt{DEC} и \texttt{FIXED} в позициях
\texttt{FORMT[2]} и \texttt{FORMT[3]} переменной присваивается тип «\texttt{D}»
(packed decimal). Также исправлена обработка случая отсутствия инициализатора:

\begin{lstlisting}[style=Cstyle,caption={Модификация подпрограммы \texttt{ODC1()}},label=lst:odc1]
int ODC1 ()
{
  /* ... */
  if ( !strcmp(FORMT[2], "BIN") && !strcmp(FORMT[3], "FIXED") )
   { SYM[ISYM].TYPE = 'B'; goto ODC11; }
  else if ( !strcmp(FORMT[2], "DEC") && !strcmp(FORMT[3], "FIXED") )
   { SYM[ISYM].TYPE = 'D'; goto ODC11; }  /* <- новая ветка */
  else
   { SYM[ISYM].TYPE = 'U'; return 2; }

ODC11:
  if ( !strcmp(FORMT[5], "INIT") )
    strcpy(SYM[ISYM++].INIT, FORMT[6]);
  else
   { SYM[ISYM].INIT[0] = '\0'; ISYM++; }  /* <- явный сброс INIT */
  return 0;
}
\end{lstlisting}

\subsection{Модификация подпрограммы \texttt{AVI2()}}

Подпрограмма второго прохода \texttt{AVI2()} была подвергнута наиболее значительной
переработке. Рассмотрим её ключевые ветви по отдельности.

\subsubsection*{Ветвь обработки одиночного операнда (\texttt{IFORMT == 1})}

\begin{lstlisting}[style=Cstyle,caption={Ветвь \texttt{AVI2()} для одиночного операнда},label=lst:avi2single]
if ( IFORMT == 1 )
 {
  if ( SYM[i].TYPE == 'B' )
   {
    /* BIN FIXED: загрузка в регистр @R2 */
    memcpy(ASS_CARD._BUFCARD.OPERAC, "LH", 2);
    strcpy(ASS_CARD._BUFCARD.OPERAND, "@R2,");
    strcat(ASS_CARD._BUFCARD.OPERAND, FORMT[0]);
    SETCOMM("Load A ->@R2");
    ZKARD();
    strcpy(RESULT_REG, "@R2");
    return 0;
   }
  else if ( SYM[i].TYPE == 'D' )
   {
    /* DEC FIXED: копирование в рабочий буфер @BUF через ZAP */
    memcpy(ASS_CARD._BUFCARD.OPERAC, "ZAP", 3);
    sprintf(opbuf, "@BUF(8),%s(%s)", FORMT[0], SYM[i].RAZR);
    memcpy(ASS_CARD.BUFCARD + 15, opbuf, strlen(opbuf));
    SETCOMM("Copy pdec B->BUF");
    ZKARD();
    return 0;
   }
  else return 3;
 }
\end{lstlisting}

\subsubsection*{Ветвь операции умножения «\texttt{*}»}

\begin{lstlisting}[style=Cstyle,caption={Ветвь \texttt{AVI2()} для операции «*»},label=lst:avi2mul]
if ( op_sign == '*' )
 {
  char const_val[10];
  char clbl[9];
  strcpy(const_val, FORMT[IFORMT-1]);
  const_label(clbl, const_val);        /* @FIVE для 5 */
  sprintf(opbuf, "@BUF(8),%s(1)", clbl);

  memcpy(ASS_CARD._BUFCARD.OPERAC, "MP", 2);
  memcpy(ASS_CARD.BUFCARD + 15, opbuf, strlen(opbuf));
  SETCOMM("Mul BUF by 5 ->BUF");
  ZKARD();

  memcpy(ASS_CARD._BUFCARD.OPERAC, "CVB", 3);
  strcpy(ASS_CARD._BUFCARD.OPERAND, "@R1,@BUF");
  SETCOMM("pdec BUF ->int @R1");
  ZKARD();
  strcpy(RESULT_REG, "@R1");

  /* регистрация константы в EXTRA[] для секции данных */
  if ( li == NEXTRA ) {
    const_label(EXTRA[NEXTRA].name, const_val);
    EXTRA[NEXTRA].type  = 'P';
    EXTRA[NEXTRA].size  = 1;
    strcpy(EXTRA[NEXTRA].value, const_val);
    NEXTRA++;
  }
  return 0;
 }
\end{lstlisting}

\subsubsection*{Ветвь оператора сравнения «\texttt{=}»}

\begin{lstlisting}[style=Cstyle,caption={Ветвь \texttt{AVI2()} для операции «=»},label=lst:avi2eq]
else if ( op_sign == '=' )
 {
  /* Сравнение: @R2 (загружен ранее) с правым операндом */
  memcpy(ASS_CARD._BUFCARD.OPERAC, "CH", 2);
  strcpy(ASS_CARD._BUFCARD.OPERAND, "@R2,");
  strcat(ASS_CARD._BUFCARD.OPERAND, FORMT[IFORMT-1]);
  SETCOMM("Compare @R2 with C");
  ZKARD();

  /* Ветвь TRUE */
  memcpy(ASS_CARD._BUFCARD.OPERAC, "BC", 2);
  strcpy(ASS_CARD._BUFCARD.OPERAND, "8,@TRUE");
  SETCOMM("Branch if equal");
  ZKARD();

  /* Ветвь FALSE */
  memcpy(ASS_CARD._BUFCARD.OPERAC, "LA", 2);
  strcpy(ASS_CARD._BUFCARD.OPERAND, "@R3,0");
  SETCOMM("Set @R3=0 (F)");
  ZKARD();

  memcpy(ASS_CARD._BUFCARD.OPERAC, "B", 1);
  strcpy(ASS_CARD._BUFCARD.OPERAND, "@SAVE");
  SETCOMM("Jump to save");
  ZKARD();

  /* Метка @TRUE, присвоение 1 */
  memcpy(ASS_CARD._BUFCARD.METKA, "@TRUE", 5);
  memcpy(ASS_CARD._BUFCARD.OPERAC, "LA", 2);
  strcpy(ASS_CARD._BUFCARD.OPERAND, "@R3,1");
  SETCOMM("Set @R3=1 (T)");
  ZKARD();

  strcpy(RESULT_REG,    "@R3");
  strcpy(PENDING_LABEL, "@SAVE");  /* метка для STH в OPA2 */
  return 0;
 }
\end{lstlisting}

\subsection{Модификация подпрограммы \texttt{OPA2()}}

В подпрограмму \texttt{OPA2()} добавлен механизм обработки отложенной метки
\texttt{PENDING\_LABEL} и использование \texttt{RESULT\_REG} для выбора
регистра-источника при инструкции \texttt{STH}:

\begin{lstlisting}[style=Cstyle,caption={Модификация подпрограммы \texttt{OPA2()}},label=lst:opa2]
int OPA2 ()
{
  /* ... */
  if ( SYM[i].TYPE == 'B' )
   {
    /* Навешиваем отложенную метку @SAVE, если она есть */
    if ( PENDING_LABEL[0] != '\0' )
     {
      strcpy(ASS_CARD._BUFCARD.METKA, PENDING_LABEL);
      ASS_CARD._BUFCARD.METKA[strlen(PENDING_LABEL)] = ' ';
      PENDING_LABEL[0] = '\0';
     }
    memcpy(ASS_CARD._BUFCARD.OPERAC, "STH", 3);
    strcpy(ASS_CARD._BUFCARD.OPERAND, RESULT_REG);
    strcat(ASS_CARD._BUFCARD.OPERAND, ",");
    strcat(ASS_CARD._BUFCARD.OPERAND, FORMT[0]);
    /* Пояснение зависит от типа результата */
    if ( !strcmp(RESULT_REG, "@R3") )
      SETCOMM("Store bool ->C");
    else
      SETCOMM("Store int ->C");
    ZKARD();
    return 0;
   }
  /* ... */
}
\end{lstlisting}

\subsection{Модификация подпрограммы \texttt{OEN2()}}

Подпрограмма эпилога \texttt{OEN2()} полностью переработана для корректной
генерации секции данных. Ключевые изменения:

\begin{lstlisting}[style=Cstyle,caption={Ключевые фрагменты модифицированной \texttt{OEN2()}},label=lst:oen2]
int OEN2 ()
{
  /* Возврат из программы */
  memcpy(ASS_CARD._BUFCARD.OPERAC, "BCR", 3);
  memcpy(ASS_CARD._BUFCARD.OPERAND, "15,@RVIX", 8);
  SETCOMM("Return to caller");
  ZKARD();

  /* Секция данных: переменные из таблицы SYM */
  for ( i = 0; i < ISYM; i++ )
   {
    if ( !isalpha(SYM[i].NAME[0]) ) continue;
    if ( SYM[i].TYPE != 'B' && SYM[i].TYPE != 'D' ) continue;
    /* ... установка метки ... */

    if ( SYM[i].TYPE == 'B' )
     {
      if ( SYM[i].INIT[0] == '\0' )
       { memcpy(ASS_CARD._BUFCARD.OPERAC, "DS", 2);
         strcpy(ASS_CARD._BUFCARD.OPERAND, "H");
         SETCOMM("Reserve halfword"); }
      else
       { memcpy(ASS_CARD._BUFCARD.OPERAC, "DC", 2);
         sprintf(plop, "H'%ld'", VALUE(SYM[i].INIT));
         strcpy(ASS_CARD._BUFCARD.OPERAND, plop);
         SETCOMM("Store halfword init"); }
     }
    else if ( SYM[i].TYPE == 'D' )
     {
      memcpy(ASS_CARD._BUFCARD.OPERAC, "DC", 2);
      sprintf(plop, "PL%s'%ld'", SYM[i].RAZR, VALUE(SYM[i].INIT));
      strcpy(ASS_CARD._BUFCARD.OPERAND, plop);
      SETCOMM("Store pdec init");
     }
    ZKARD();
   }

  /* Дополнительные константы (из EXTRA[]) */
  for ( ei = 0; ei < NEXTRA; ei++ )
   {
    strcpy(ASS_CARD._BUFCARD.METKA, EXTRA[ei].name);
    memcpy(ASS_CARD._BUFCARD.OPERAC, "DC", 2);
    sprintf(plop, "PL%d'%s'", EXTRA[ei].size, EXTRA[ei].value);
    strcpy(ASS_CARD._BUFCARD.OPERAND, plop);
    SETCOMM("Store packed const");
    ZKARD();
   }

  /* Буфер @BUF для упакованной десятичной арифметики */
  if ( NEXTRA > 0 )
   {
    memcpy(ASS_CARD._BUFCARD.OPERAC, "DS", 2);
    strcpy(ASS_CARD._BUFCARD.OPERAND, "0F");
    SETCOMM("Align next item"); ZKARD();

    memcpy(ASS_CARD._BUFCARD.METKA, "@BUF", 4);
    memcpy(ASS_CARD._BUFCARD.OPERAC, "DC", 2);
    strcpy(ASS_CARD._BUFCARD.OPERAND, "PL8'0'");
    SETCOMM("Clear packed BUF"); ZKARD();
   }

  /* Определения регистров через EQU */
  /* @RBASE=5, @R1=6, @R2=2, @R3=3, @RVIX=14 */

  /* Завершение блока */
  memcpy(ASS_CARD._BUFCARD.OPERAC, "END", 3);
  SETCOMM("Program ends");
  ZKARD();
  return 0;
}
\end{lstlisting}

\subsection{Модификация подпрограммы \texttt{OPR2()}}

В подпрограмме пролога \texttt{OPR2()} изменены имена регистров
(добавлен префикс «\texttt{@}»), а также заменены длинные жёстко закодированные
комментарии на вызовы \texttt{SETCOMM()}:

\begin{lstlisting}[style=Cstyle,caption={Модификация подпрограммы \texttt{OPR2()}},label=lst:opr2]
int OPR2 ()
{
  /* START */
  /* имя программы -> поле METKA */
  memcpy(ASS_CARD._BUFCARD.OPERAC, "START", 5);
  memcpy(ASS_CARD._BUFCARD.OPERAND, "0", 1);
  SETCOMM("Init rel addr=0");
  ZKARD();

  /* BALR */
  memcpy(ASS_CARD._BUFCARD.OPERAC, "BALR", 4);
  memcpy(ASS_CARD._BUFCARD.OPERAND, "@RBASE,0", 8);
  ASS_CARD._BUFCARD.OPERAND[8] = ' ';
  SETCOMM("Set @RBASE=cur PC");
  ZKARD();

  /* USING */
  memcpy(ASS_CARD._BUFCARD.OPERAC, "USING", 5);
  memcpy(ASS_CARD._BUFCARD.OPERAND, "*,@RBASE", 8);
  ASS_CARD._BUFCARD.OPERAND[8] = ' ';
  SETCOMM("Use @RBASE for addr");
  ZKARD();
  return 0;
}
\end{lstlisting}

% ============================================================
% РАЗДЕЛ 5. РЕЗУЛЬТАТЫ ЭКСПЕРИМЕНТАЛЬНОЙ ПРОВЕРКИ
% ============================================================
\section{Результаты экспериментальной проверки модифицированного компилятора с ЯВУ}

Работоспособность модифицированного макетного компилятора с~ЯВУ подтверждена
испытательным прогоном в среде ОС~Linux. Компиляция выполнялась в~32-битном
режиме внутри Docker-контейнера на базе образа Ubuntu~22.04 с использованием
компилятора \texttt{gcc~-m32}.

Команды сборки и запуска:

\begin{lstlisting}[style=ASMstyle,caption={Команды сборки и запуска компилятора},label=lst:build]
docker build -t lab32 .
docker run -it --rm -v "$PWD"/src:/lab -w /lab lab32
make
./build/komppl.exe step1/ex11.pli
\end{lstlisting}

Результатом компиляции является файл \texttt{ex11.ass}. Его содержимое совпало с
эталоном (рис.~\ref{fig:etalon}) и приведено в листинге~\ref{lst:result_asm}.

\begin{figure}[H]
\begin{lstlisting}[style=ASMstyle,caption={Полученный ассемблерный эквивалент \texttt{ex11.ass}},label=lst:result_asm]
EX11     START 0            Init rel addr=0
         BALR  @RBASE,0     Set @RBASE=cur PC
         USING *,@RBASE     Use @RBASE for addr
         ZAP   @BUF(8),B(3) Copy pdec B->BUF
         MP    @BUF(8),@FIVE(1) Mul BUF by 5 ->BUF
         CVB   @R1,@BUF     pdec BUF ->int @R1
         STH   @R1,C        Store int ->C
         LH    @R2,A        Load A ->@R2
         CH    @R2,C        Compare @R2 with C
         BC    8,@TRUE      Branch if equal
         LA    @R3,0        Set @R3=0 (F)
         B     @SAVE        Jump to save
@TRUE    LA    @R3,1        Set @R3=1 (T)
@SAVE    STH   @R3,C        Store bool ->C
         BCR   15,@RVIX     Return to caller
A        DC    H'3'         Store halfword init
B        DC    PL3'3'       Store pdec init
C        DS    H            Reserve halfword
@FIVE    DC    PL1'5'       Store packed const
         DS    0F           Align next item
@BUF     DC    PL8'0'       Clear packed BUF
@RBASE   EQU   5            Set @RBASE=5
@R1      EQU   6            Set @R1=6
@R2      EQU   2            Set @R2=2
@R3      EQU   3            Set @R3=3
@RVIX    EQU   14           Set @RVIX=14
         END                Program ends
\end{lstlisting}
\end{figure}

\subsection*{Пояснения к сгенерированному тексту на Ассемблере}

\textbf{Семантика исходной программы:}
\begin{itemize}
  \item Литерал \texttt{11B} = $1_2 \cdot 2^1 + 1_2 \cdot 2^0 = 3_{10}$, следовательно:
        $\texttt{A} = 3$, $\texttt{B} = 3$;
  \item \texttt{C = B * 5}: упакованное десятичное $B=3$ умножается на~5,
        результат $15$ конвертируется в двоичное целое и сохраняется в \texttt{C};
  \item \texttt{C = A = C}: сравнение $A=3$ и $C=15$; результат~--- \textit{ложь},
        поэтому $\texttt{C} = 0$.
\end{itemize}

\textbf{Пролог} (строки \texttt{START}~– \texttt{USING}): сформирован подпрограммой
\texttt{OPR2()}; устанавливает нулевой начальный адрес, загружает базовый регистр
\texttt{@RBASE} (=~регистр~5) и объявляет его в директиве \texttt{USING}.

\textbf{Оператор \texttt{C = B * 5}}:
\begin{itemize}
  \item \texttt{ZAP @BUF(8),B(3)} — копирование упакованного десятичного~$B$ в
        8-байтный рабочий буфер \texttt{@BUF};
  \item \texttt{MP @BUF(8),@FIVE(1)} — умножение содержимого \texttt{@BUF} на
        упакованную константу \texttt{@FIVE} (=~$5$);
  \item \texttt{CVB @R1,@BUF} — конвертация результата в двоичное целое,
        размещение в \texttt{@R1} (= регистр~6);
  \item \texttt{STH @R1,C} — запись halfword из \texttt{@R1} в ячейку \texttt{C}.
\end{itemize}

\textbf{Оператор \texttt{C = A = C}}:
\begin{itemize}
  \item \texttt{LH @R2,A} — загрузка переменной~$A$ в регистр \texttt{@R2} (= рег.~2);
  \item \texttt{CH @R2,C} — сравнение \texttt{@R2} с \texttt{C};
  \item \texttt{BC 8,@TRUE} — ветвление при равенстве (маска~8);
  \item ветвь \textit{ложь}: \texttt{LA @R3,0}; \texttt{B @SAVE};
  \item ветвь \textit{истина}: \texttt{@TRUE: LA @R3,1};
  \item \texttt{@SAVE: STH @R3,C} — запись булева результата в \texttt{C}
        (метка \texttt{@SAVE} навешена механизмом \texttt{PENDING\_LABEL}).
\end{itemize}

\textbf{Эпилог} (строки \texttt{BCR}~– \texttt{END}): сформирован подпрограммой
\texttt{OEN2()}; включает инструкцию возврата, секцию данных (переменные
\texttt{A}, \texttt{B}, \texttt{C}), константу \texttt{@FIVE}, рабочий буфер
\texttt{@BUF} с выравниванием \texttt{DS 0F}, определения регистров через
\texttt{EQU} и директиву \texttt{END}.

% ============================================================
% РАЗДЕЛ 6. ВЫВОДЫ
% ============================================================
\section{Выводы по результатам работы}

Порученное задание выполнено успешно. Успешность выполнения порученного задания
подтверждена полным совпадением выходного текста, полученного при
экспериментальной проверке компилятора с~ЯВУ, с эталоном на рис.~\ref{fig:etalon}.

В ходе выполнения работы:

\begin{enumerate}
  \item Произведена модификация базы данных лексического анализатора: в перечень
        терминальных символов-разделителей добавлен символ умножения «\texttt{*}».

  \item Произведена модификация базы данных синтаксического анализатора:
        таблица \texttt{SINT[]} расширена с~201 до~232 записей; добавлены
        ветви для объявления переменных типа \texttt{DEC FIXED}, операции
        умножения «\texttt{*}», операции сравнения «\texttt{=}» и использования
        целочисленного литерала \texttt{<RZR>} в качестве операнда выражения;
        соответствующие изменения внесены в таблицу входов \texttt{VXOD[]} и
        матрицу предшествования \texttt{TPR[][]}.

  \item Произведена модификация семантической библиотеки компилятора:
        подпрограммы \texttt{ODC1()}, \texttt{AVI2()}, \texttt{OPA2()},
        \texttt{OEN2()}, \texttt{OPR2()} изменены в соответствии с требованиями
        персонального варианта; введены вспомогательные подпрограммы
        \texttt{SETCOMM()} и \texttt{const\_label()}, а также новые глобальные
        переменные для передачи состояния между семантическими подпрограммами.

  \item Работоспособность модифицированного компилятора подтверждена
        испытательным прогоном на тестовом примере \texttt{ex11.pli} в среде
        ОС~Linux.
\end{enumerate}

К позитивным аспектам выполненной работы следует отнести следующие:
архитектура макета компилятора (параметрическая настройка через таблицы БД и
библиотеку семантических подпрограмм) обеспечила относительно локальный характер
всех необходимых изменений и позволила свести к минимуму риск нежелательных
побочных эффектов.

\end{document}
